

Bonus\\
\vspace{.4cm}
1. Yes, the PCA and KPCA reutrn the same results if KPCA uses a linear kernel. See plots below.\\
\vspace{.4cm}
2. Let V be of maximaal dimension d and compare the complexity of each algorithm.
\vspace{.4cm}
The PCA is an optimization problem with cost function $\sum\limits_{j=1}^d v_j (frac{1}{n}sum\limits_{i=1}^n x_i x_i^{\top})$,\\
where the the optimation problem for KPCA uses $frac{1}{n}\sum\limits{j=1}^d \alpha_j^{top} K^2 \alpha_j$.\\
\vspace{.4cm}
Upon comparison of the two cost functions the difference is complexity is mainly given by the difference between
calculating $sum\limits_{i=1}^n x_i x_i^{\top}$ vs. calculating K^2. Let's analyse the number of operation for each.
\vspace{.4cm}
$x_i$ is has d rows so $x_i x_^{\top}$ requires d\timesd products and the sum over n requires n\timees d^2 products.
$For each entry in $K^2$ there are n products and n nummations. Since there are $n^2$ matrix entries there are n^3 products\\
and $n^3$ summations.
\vspace{.4cm}
Since products are much more compuational expensive the difference between thee two algoithms is between the calaculation of\\
$n \times d^2$ and $n^3 terms. Since d is general is less than n, PCA will generally be faster than KPCA for a linear kernel.\\
only then the dimension of V is large compared with the number of samples would KPCA be a better choice in this case.
Using the data problem 3 with 275 samples with d=2, the elapsed time for PCA was .5 ms and for KPCA 3 ms.
\include{STAT672Homework1_bonus.pdf}
